\documentclass[11pt]{article}
% =======================================================================
%   Shared style for CS 300 notes.
%   Usage: place `% =======================================================================
%   Shared style for CS 300 notes.
%   Usage: place `% =======================================================================
%   Shared style for CS 300 notes.
%   Usage: place `\input{../../notes-style.tex}` right after \documentclass
%   in each chapter's lectures.tex and notes.tex.
% =======================================================================

% ---- Layout -----------------------------------------------------------
\usepackage[margin=1in]{geometry}
\usepackage{parskip}        % space between paragraphs, no indent
\usepackage{microtype}      % subtle kerning/spacing improvements

% ---- Colors (muted palette, easy on light-sensitive eyes) -------------
\usepackage{xcolor}
\definecolor{pagebg}       {HTML}{FAF6F0}  % warm cream background
\definecolor{bodyfg}       {HTML}{3D3A36}  % warm dark gray body text
\definecolor{heading}      {HTML}{3D6A6B}  % muted teal
\definecolor{subheading}   {HTML}{5B7A9C}  % dusty blue
\definecolor{accent}       {HTML}{8A6B4A}  % warm tan
\definecolor{linkcol}      {HTML}{5B7A9C}  % same as subheading
\definecolor{mutedtext}    {HTML}{6B6560}  % secondary / caption text
\definecolor{rulecol}      {HTML}{C8BFB0}  % separator / rule

% Callout box palette
\definecolor{defnFrame}    {HTML}{9DB89B}  \definecolor{defnBack}    {HTML}{EEF2E8}
\definecolor{ideaFrame}    {HTML}{9FB4CD}  \definecolor{ideaBack}    {HTML}{E8EEF5}
\definecolor{warnFrame}    {HTML}{CFAE87}  \definecolor{warnBack}    {HTML}{F5EDDF}
\definecolor{noteFrame}    {HTML}{BFA6AE}  \definecolor{noteBack}    {HTML}{F2E8EE}
\definecolor{exFrame}      {HTML}{9FBDBD}  \definecolor{exBack}      {HTML}{E8F0F0}
\definecolor{codeBack}     {HTML}{F3EDE2}
\definecolor{codeBorder}   {HTML}{D4CBBA}
\definecolor{codeKeyword}  {HTML}{6B4F7A}
\definecolor{codeString}   {HTML}{8A6B4A}
\definecolor{codeComment}  {HTML}{8A857F}

% ---- Page background + base text color --------------------------------
\usepackage{pagecolor}
\pagecolor{pagebg}
\color{bodyfg}

% ---- Fonts ------------------------------------------------------------
\usepackage[T1]{fontenc}
\IfFileExists{libertine.sty}{%
    \usepackage{libertine}      % warm, readable serif
    \usepackage[scaled=0.95]{inconsolata}  % softer monospace
}{}

% ---- Hyperref ---------------------------------------------------------
\usepackage{hyperref}
\hypersetup{
    colorlinks=true,
    linkcolor=linkcol,
    urlcolor=linkcol,
    citecolor=linkcol,
    pdfborder={0 0 0},
}

% ---- Section styling --------------------------------------------------
\usepackage[explicit]{titlesec}
\titleformat{\section}
    {\color{heading}\Large\bfseries}
    {\color{heading}\thesection\hspace{0.6em}}{0pt}{#1}
\titleformat{\subsection}
    {\color{subheading}\large\bfseries}
    {\color{subheading}\thesubsection\hspace{0.5em}}{0pt}{#1}
\titleformat{\subsubsection}
    {\color{accent}\normalsize\bfseries}
    {\color{accent}\thesubsubsection\hspace{0.5em}}{0pt}{#1}
\titleformat{\paragraph}[runin]
    {\color{accent}\bfseries}{}{0pt}{#1.}

\titlespacing*{\section}{0pt}{1.6em}{0.6em}
\titlespacing*{\subsection}{0pt}{1.2em}{0.4em}

% ---- Enumerate / itemize spacing --------------------------------------
\usepackage{enumitem}
\setlist[itemize]{leftmargin=1.4em, itemsep=2pt, topsep=4pt}
\setlist[enumerate]{leftmargin=1.6em, itemsep=2pt, topsep=4pt}

% ---- Code listings ----------------------------------------------------
\usepackage{listings}
\lstdefinestyle{notescode}{
    backgroundcolor=\color{codeBack},
    basicstyle=\ttfamily\small\color{bodyfg},
    keywordstyle=\color{codeKeyword}\bfseries,
    stringstyle=\color{codeString},
    commentstyle=\color{codeComment}\itshape,
    frame=single,
    framesep=6pt,
    rulecolor=\color{codeBorder},
    showstringspaces=false,
    breaklines=true,
    xleftmargin=10pt,
    xrightmargin=10pt,
    aboveskip=8pt,
    belowskip=8pt,
}
\lstset{style=notescode}

% ---- Callout boxes (tcolorbox) ----------------------------------------
\usepackage[most]{tcolorbox}

\newtcolorbox{defnbox}[1][]{
    enhanced, breakable,
    colback=defnBack, colframe=defnFrame, coltitle=bodyfg,
    title={\bfseries Definition \textemdash\ #1},
    fonttitle=\bfseries,
    boxrule=0.5pt, arc=2pt, left=8pt, right=8pt, top=6pt, bottom=6pt,
    attach boxed title to top left={xshift=10pt, yshift=-6pt},
    boxed title style={colback=defnFrame, colframe=defnFrame, arc=2pt},
    before skip=8pt, after skip=8pt,
}
\newtcolorbox{ideabox}[1][]{
    enhanced, breakable,
    colback=ideaBack, colframe=ideaFrame, coltitle=bodyfg,
    title={\bfseries Key idea #1},
    boxrule=0.5pt, arc=2pt, left=8pt, right=8pt, top=6pt, bottom=6pt,
    attach boxed title to top left={xshift=10pt, yshift=-6pt},
    boxed title style={colback=ideaFrame, colframe=ideaFrame, arc=2pt},
    before skip=8pt, after skip=8pt,
}
\newtcolorbox{warnbox}[1][]{
    enhanced, breakable,
    colback=warnBack, colframe=warnFrame, coltitle=bodyfg,
    title={\bfseries Gotcha #1},
    boxrule=0.5pt, arc=2pt, left=8pt, right=8pt, top=6pt, bottom=6pt,
    attach boxed title to top left={xshift=10pt, yshift=-6pt},
    boxed title style={colback=warnFrame, colframe=warnFrame, arc=2pt},
    before skip=8pt, after skip=8pt,
}
\newtcolorbox{notebox}[1][]{
    enhanced, breakable,
    colback=noteBack, colframe=noteFrame, coltitle=bodyfg,
    title={\bfseries Aside #1},
    boxrule=0.5pt, arc=2pt, left=8pt, right=8pt, top=6pt, bottom=6pt,
    attach boxed title to top left={xshift=10pt, yshift=-6pt},
    boxed title style={colback=noteFrame, colframe=noteFrame, arc=2pt},
    before skip=8pt, after skip=8pt,
}
\newtcolorbox{examplebox}[1][]{
    enhanced, breakable,
    colback=exBack, colframe=exFrame, coltitle=bodyfg,
    title={\bfseries Example #1},
    boxrule=0.5pt, arc=2pt, left=8pt, right=8pt, top=6pt, bottom=6pt,
    attach boxed title to top left={xshift=10pt, yshift=-6pt},
    boxed title style={colback=exFrame, colframe=exFrame, arc=2pt},
    before skip=8pt, after skip=8pt,
}

% ---- TikZ graph styles (added 2026-04-25 for ch_10 augmentation) ------
% tcolorbox[most] already loads tikz; this block declares the libraries
% + shared `vertex` / `edge` styles used by ch_10's general directed-graph
% diagrams. ch_9's tree-style TikZ uses inline overrides and is
% unaffected. Simple circle nodes + arrow edges only -- no production
% styling.
\usetikzlibrary{arrows.meta, positioning, calc}
\tikzset{
    vertex/.style={circle, draw, minimum size=7mm, inner sep=0pt,
                   font=\small, thick},
    vsource/.style={vertex, fill=ideaBack},
    vdone/.style={vertex, fill=defnBack},
    edge/.style={->, >=Stealth, thick},
    uedge/.style={-, thick},
    treeedge/.style={->, >=Stealth, thick},
    backedge/.style={->, >=Stealth, thick, dashed, red!70!black},
    fwdedge/.style={->, >=Stealth, thick, blue!60!black},
    crossedge/.style={->, >=Stealth, thick, green!50!black, dotted},
    mstedge/.style={-, very thick, accent},
}

% ---- Title block ------------------------------------------------------
\usepackage{titling}
\pretitle{\begin{center}\color{heading}\LARGE\bfseries}
\posttitle{\end{center}\vskip -0.4em}
\preauthor{\begin{center}\color{mutedtext}\normalsize}
\postauthor{\end{center}}
\predate{\begin{center}\color{mutedtext}\small}
\postdate{\end{center}\vspace{0.4em}{\color{rulecol}\hrule height 0.4pt}\vspace{1.2em}}
` right after \documentclass
%   in each chapter's lectures.tex and notes.tex.
% =======================================================================

% ---- Layout -----------------------------------------------------------
\usepackage[margin=1in]{geometry}
\usepackage{parskip}        % space between paragraphs, no indent
\usepackage{microtype}      % subtle kerning/spacing improvements

% ---- Colors (muted palette, easy on light-sensitive eyes) -------------
\usepackage{xcolor}
\definecolor{pagebg}       {HTML}{FAF6F0}  % warm cream background
\definecolor{bodyfg}       {HTML}{3D3A36}  % warm dark gray body text
\definecolor{heading}      {HTML}{3D6A6B}  % muted teal
\definecolor{subheading}   {HTML}{5B7A9C}  % dusty blue
\definecolor{accent}       {HTML}{8A6B4A}  % warm tan
\definecolor{linkcol}      {HTML}{5B7A9C}  % same as subheading
\definecolor{mutedtext}    {HTML}{6B6560}  % secondary / caption text
\definecolor{rulecol}      {HTML}{C8BFB0}  % separator / rule

% Callout box palette
\definecolor{defnFrame}    {HTML}{9DB89B}  \definecolor{defnBack}    {HTML}{EEF2E8}
\definecolor{ideaFrame}    {HTML}{9FB4CD}  \definecolor{ideaBack}    {HTML}{E8EEF5}
\definecolor{warnFrame}    {HTML}{CFAE87}  \definecolor{warnBack}    {HTML}{F5EDDF}
\definecolor{noteFrame}    {HTML}{BFA6AE}  \definecolor{noteBack}    {HTML}{F2E8EE}
\definecolor{exFrame}      {HTML}{9FBDBD}  \definecolor{exBack}      {HTML}{E8F0F0}
\definecolor{codeBack}     {HTML}{F3EDE2}
\definecolor{codeBorder}   {HTML}{D4CBBA}
\definecolor{codeKeyword}  {HTML}{6B4F7A}
\definecolor{codeString}   {HTML}{8A6B4A}
\definecolor{codeComment}  {HTML}{8A857F}

% ---- Page background + base text color --------------------------------
\usepackage{pagecolor}
\pagecolor{pagebg}
\color{bodyfg}

% ---- Fonts ------------------------------------------------------------
\usepackage[T1]{fontenc}
\IfFileExists{libertine.sty}{%
    \usepackage{libertine}      % warm, readable serif
    \usepackage[scaled=0.95]{inconsolata}  % softer monospace
}{}

% ---- Hyperref ---------------------------------------------------------
\usepackage{hyperref}
\hypersetup{
    colorlinks=true,
    linkcolor=linkcol,
    urlcolor=linkcol,
    citecolor=linkcol,
    pdfborder={0 0 0},
}

% ---- Section styling --------------------------------------------------
\usepackage[explicit]{titlesec}
\titleformat{\section}
    {\color{heading}\Large\bfseries}
    {\color{heading}\thesection\hspace{0.6em}}{0pt}{#1}
\titleformat{\subsection}
    {\color{subheading}\large\bfseries}
    {\color{subheading}\thesubsection\hspace{0.5em}}{0pt}{#1}
\titleformat{\subsubsection}
    {\color{accent}\normalsize\bfseries}
    {\color{accent}\thesubsubsection\hspace{0.5em}}{0pt}{#1}
\titleformat{\paragraph}[runin]
    {\color{accent}\bfseries}{}{0pt}{#1.}

\titlespacing*{\section}{0pt}{1.6em}{0.6em}
\titlespacing*{\subsection}{0pt}{1.2em}{0.4em}

% ---- Enumerate / itemize spacing --------------------------------------
\usepackage{enumitem}
\setlist[itemize]{leftmargin=1.4em, itemsep=2pt, topsep=4pt}
\setlist[enumerate]{leftmargin=1.6em, itemsep=2pt, topsep=4pt}

% ---- Code listings ----------------------------------------------------
\usepackage{listings}
\lstdefinestyle{notescode}{
    backgroundcolor=\color{codeBack},
    basicstyle=\ttfamily\small\color{bodyfg},
    keywordstyle=\color{codeKeyword}\bfseries,
    stringstyle=\color{codeString},
    commentstyle=\color{codeComment}\itshape,
    frame=single,
    framesep=6pt,
    rulecolor=\color{codeBorder},
    showstringspaces=false,
    breaklines=true,
    xleftmargin=10pt,
    xrightmargin=10pt,
    aboveskip=8pt,
    belowskip=8pt,
}
\lstset{style=notescode}

% ---- Callout boxes (tcolorbox) ----------------------------------------
\usepackage[most]{tcolorbox}

\newtcolorbox{defnbox}[1][]{
    enhanced, breakable,
    colback=defnBack, colframe=defnFrame, coltitle=bodyfg,
    title={\bfseries Definition \textemdash\ #1},
    fonttitle=\bfseries,
    boxrule=0.5pt, arc=2pt, left=8pt, right=8pt, top=6pt, bottom=6pt,
    attach boxed title to top left={xshift=10pt, yshift=-6pt},
    boxed title style={colback=defnFrame, colframe=defnFrame, arc=2pt},
    before skip=8pt, after skip=8pt,
}
\newtcolorbox{ideabox}[1][]{
    enhanced, breakable,
    colback=ideaBack, colframe=ideaFrame, coltitle=bodyfg,
    title={\bfseries Key idea #1},
    boxrule=0.5pt, arc=2pt, left=8pt, right=8pt, top=6pt, bottom=6pt,
    attach boxed title to top left={xshift=10pt, yshift=-6pt},
    boxed title style={colback=ideaFrame, colframe=ideaFrame, arc=2pt},
    before skip=8pt, after skip=8pt,
}
\newtcolorbox{warnbox}[1][]{
    enhanced, breakable,
    colback=warnBack, colframe=warnFrame, coltitle=bodyfg,
    title={\bfseries Gotcha #1},
    boxrule=0.5pt, arc=2pt, left=8pt, right=8pt, top=6pt, bottom=6pt,
    attach boxed title to top left={xshift=10pt, yshift=-6pt},
    boxed title style={colback=warnFrame, colframe=warnFrame, arc=2pt},
    before skip=8pt, after skip=8pt,
}
\newtcolorbox{notebox}[1][]{
    enhanced, breakable,
    colback=noteBack, colframe=noteFrame, coltitle=bodyfg,
    title={\bfseries Aside #1},
    boxrule=0.5pt, arc=2pt, left=8pt, right=8pt, top=6pt, bottom=6pt,
    attach boxed title to top left={xshift=10pt, yshift=-6pt},
    boxed title style={colback=noteFrame, colframe=noteFrame, arc=2pt},
    before skip=8pt, after skip=8pt,
}
\newtcolorbox{examplebox}[1][]{
    enhanced, breakable,
    colback=exBack, colframe=exFrame, coltitle=bodyfg,
    title={\bfseries Example #1},
    boxrule=0.5pt, arc=2pt, left=8pt, right=8pt, top=6pt, bottom=6pt,
    attach boxed title to top left={xshift=10pt, yshift=-6pt},
    boxed title style={colback=exFrame, colframe=exFrame, arc=2pt},
    before skip=8pt, after skip=8pt,
}

% ---- TikZ graph styles (added 2026-04-25 for ch_10 augmentation) ------
% tcolorbox[most] already loads tikz; this block declares the libraries
% + shared `vertex` / `edge` styles used by ch_10's general directed-graph
% diagrams. ch_9's tree-style TikZ uses inline overrides and is
% unaffected. Simple circle nodes + arrow edges only -- no production
% styling.
\usetikzlibrary{arrows.meta, positioning, calc}
\tikzset{
    vertex/.style={circle, draw, minimum size=7mm, inner sep=0pt,
                   font=\small, thick},
    vsource/.style={vertex, fill=ideaBack},
    vdone/.style={vertex, fill=defnBack},
    edge/.style={->, >=Stealth, thick},
    uedge/.style={-, thick},
    treeedge/.style={->, >=Stealth, thick},
    backedge/.style={->, >=Stealth, thick, dashed, red!70!black},
    fwdedge/.style={->, >=Stealth, thick, blue!60!black},
    crossedge/.style={->, >=Stealth, thick, green!50!black, dotted},
    mstedge/.style={-, very thick, accent},
}

% ---- Title block ------------------------------------------------------
\usepackage{titling}
\pretitle{\begin{center}\color{heading}\LARGE\bfseries}
\posttitle{\end{center}\vskip -0.4em}
\preauthor{\begin{center}\color{mutedtext}\normalsize}
\postauthor{\end{center}}
\predate{\begin{center}\color{mutedtext}\small}
\postdate{\end{center}\vspace{0.4em}{\color{rulecol}\hrule height 0.4pt}\vspace{1.2em}}
` right after \documentclass
%   in each chapter's lectures.tex and notes.tex.
% =======================================================================

% ---- Layout -----------------------------------------------------------
\usepackage[margin=1in]{geometry}
\usepackage{parskip}        % space between paragraphs, no indent
\usepackage{microtype}      % subtle kerning/spacing improvements

% ---- Colors (muted palette, easy on light-sensitive eyes) -------------
\usepackage{xcolor}
\definecolor{pagebg}       {HTML}{FAF6F0}  % warm cream background
\definecolor{bodyfg}       {HTML}{3D3A36}  % warm dark gray body text
\definecolor{heading}      {HTML}{3D6A6B}  % muted teal
\definecolor{subheading}   {HTML}{5B7A9C}  % dusty blue
\definecolor{accent}       {HTML}{8A6B4A}  % warm tan
\definecolor{linkcol}      {HTML}{5B7A9C}  % same as subheading
\definecolor{mutedtext}    {HTML}{6B6560}  % secondary / caption text
\definecolor{rulecol}      {HTML}{C8BFB0}  % separator / rule

% Callout box palette
\definecolor{defnFrame}    {HTML}{9DB89B}  \definecolor{defnBack}    {HTML}{EEF2E8}
\definecolor{ideaFrame}    {HTML}{9FB4CD}  \definecolor{ideaBack}    {HTML}{E8EEF5}
\definecolor{warnFrame}    {HTML}{CFAE87}  \definecolor{warnBack}    {HTML}{F5EDDF}
\definecolor{noteFrame}    {HTML}{BFA6AE}  \definecolor{noteBack}    {HTML}{F2E8EE}
\definecolor{exFrame}      {HTML}{9FBDBD}  \definecolor{exBack}      {HTML}{E8F0F0}
\definecolor{codeBack}     {HTML}{F3EDE2}
\definecolor{codeBorder}   {HTML}{D4CBBA}
\definecolor{codeKeyword}  {HTML}{6B4F7A}
\definecolor{codeString}   {HTML}{8A6B4A}
\definecolor{codeComment}  {HTML}{8A857F}

% ---- Page background + base text color --------------------------------
\usepackage{pagecolor}
\pagecolor{pagebg}
\color{bodyfg}

% ---- Fonts ------------------------------------------------------------
\usepackage[T1]{fontenc}
\IfFileExists{libertine.sty}{%
    \usepackage{libertine}      % warm, readable serif
    \usepackage[scaled=0.95]{inconsolata}  % softer monospace
}{}

% ---- Hyperref ---------------------------------------------------------
\usepackage{hyperref}
\hypersetup{
    colorlinks=true,
    linkcolor=linkcol,
    urlcolor=linkcol,
    citecolor=linkcol,
    pdfborder={0 0 0},
}

% ---- Section styling --------------------------------------------------
\usepackage[explicit]{titlesec}
\titleformat{\section}
    {\color{heading}\Large\bfseries}
    {\color{heading}\thesection\hspace{0.6em}}{0pt}{#1}
\titleformat{\subsection}
    {\color{subheading}\large\bfseries}
    {\color{subheading}\thesubsection\hspace{0.5em}}{0pt}{#1}
\titleformat{\subsubsection}
    {\color{accent}\normalsize\bfseries}
    {\color{accent}\thesubsubsection\hspace{0.5em}}{0pt}{#1}
\titleformat{\paragraph}[runin]
    {\color{accent}\bfseries}{}{0pt}{#1.}

\titlespacing*{\section}{0pt}{1.6em}{0.6em}
\titlespacing*{\subsection}{0pt}{1.2em}{0.4em}

% ---- Enumerate / itemize spacing --------------------------------------
\usepackage{enumitem}
\setlist[itemize]{leftmargin=1.4em, itemsep=2pt, topsep=4pt}
\setlist[enumerate]{leftmargin=1.6em, itemsep=2pt, topsep=4pt}

% ---- Code listings ----------------------------------------------------
\usepackage{listings}
\lstdefinestyle{notescode}{
    backgroundcolor=\color{codeBack},
    basicstyle=\ttfamily\small\color{bodyfg},
    keywordstyle=\color{codeKeyword}\bfseries,
    stringstyle=\color{codeString},
    commentstyle=\color{codeComment}\itshape,
    frame=single,
    framesep=6pt,
    rulecolor=\color{codeBorder},
    showstringspaces=false,
    breaklines=true,
    xleftmargin=10pt,
    xrightmargin=10pt,
    aboveskip=8pt,
    belowskip=8pt,
}
\lstset{style=notescode}

% ---- Callout boxes (tcolorbox) ----------------------------------------
\usepackage[most]{tcolorbox}

\newtcolorbox{defnbox}[1][]{
    enhanced, breakable,
    colback=defnBack, colframe=defnFrame, coltitle=bodyfg,
    title={\bfseries Definition \textemdash\ #1},
    fonttitle=\bfseries,
    boxrule=0.5pt, arc=2pt, left=8pt, right=8pt, top=6pt, bottom=6pt,
    attach boxed title to top left={xshift=10pt, yshift=-6pt},
    boxed title style={colback=defnFrame, colframe=defnFrame, arc=2pt},
    before skip=8pt, after skip=8pt,
}
\newtcolorbox{ideabox}[1][]{
    enhanced, breakable,
    colback=ideaBack, colframe=ideaFrame, coltitle=bodyfg,
    title={\bfseries Key idea #1},
    boxrule=0.5pt, arc=2pt, left=8pt, right=8pt, top=6pt, bottom=6pt,
    attach boxed title to top left={xshift=10pt, yshift=-6pt},
    boxed title style={colback=ideaFrame, colframe=ideaFrame, arc=2pt},
    before skip=8pt, after skip=8pt,
}
\newtcolorbox{warnbox}[1][]{
    enhanced, breakable,
    colback=warnBack, colframe=warnFrame, coltitle=bodyfg,
    title={\bfseries Gotcha #1},
    boxrule=0.5pt, arc=2pt, left=8pt, right=8pt, top=6pt, bottom=6pt,
    attach boxed title to top left={xshift=10pt, yshift=-6pt},
    boxed title style={colback=warnFrame, colframe=warnFrame, arc=2pt},
    before skip=8pt, after skip=8pt,
}
\newtcolorbox{notebox}[1][]{
    enhanced, breakable,
    colback=noteBack, colframe=noteFrame, coltitle=bodyfg,
    title={\bfseries Aside #1},
    boxrule=0.5pt, arc=2pt, left=8pt, right=8pt, top=6pt, bottom=6pt,
    attach boxed title to top left={xshift=10pt, yshift=-6pt},
    boxed title style={colback=noteFrame, colframe=noteFrame, arc=2pt},
    before skip=8pt, after skip=8pt,
}
\newtcolorbox{examplebox}[1][]{
    enhanced, breakable,
    colback=exBack, colframe=exFrame, coltitle=bodyfg,
    title={\bfseries Example #1},
    boxrule=0.5pt, arc=2pt, left=8pt, right=8pt, top=6pt, bottom=6pt,
    attach boxed title to top left={xshift=10pt, yshift=-6pt},
    boxed title style={colback=exFrame, colframe=exFrame, arc=2pt},
    before skip=8pt, after skip=8pt,
}

% ---- TikZ graph styles (added 2026-04-25 for ch_10 augmentation) ------
% tcolorbox[most] already loads tikz; this block declares the libraries
% + shared `vertex` / `edge` styles used by ch_10's general directed-graph
% diagrams. ch_9's tree-style TikZ uses inline overrides and is
% unaffected. Simple circle nodes + arrow edges only -- no production
% styling.
\usetikzlibrary{arrows.meta, positioning, calc}
\tikzset{
    vertex/.style={circle, draw, minimum size=7mm, inner sep=0pt,
                   font=\small, thick},
    vsource/.style={vertex, fill=ideaBack},
    vdone/.style={vertex, fill=defnBack},
    edge/.style={->, >=Stealth, thick},
    uedge/.style={-, thick},
    treeedge/.style={->, >=Stealth, thick},
    backedge/.style={->, >=Stealth, thick, dashed, red!70!black},
    fwdedge/.style={->, >=Stealth, thick, blue!60!black},
    crossedge/.style={->, >=Stealth, thick, green!50!black, dotted},
    mstedge/.style={-, very thick, accent},
}

% ---- Title block ------------------------------------------------------
\usepackage{titling}
\pretitle{\begin{center}\color{heading}\LARGE\bfseries}
\posttitle{\end{center}\vskip -0.4em}
\preauthor{\begin{center}\color{mutedtext}\normalsize}
\postauthor{\end{center}}
\predate{\begin{center}\color{mutedtext}\small}
\postdate{\end{center}\vspace{0.4em}{\color{rulecol}\hrule height 0.4pt}\vspace{1.2em}}


\title{CS 300 -- Chapter 5 Lectures\\\large Hash Tables and Cryptographic Hashing}
\author{Jose de Lima}
\date{\today}

\begin{document}
\maketitle

\begin{ideabox}[Chapter map]
\textbf{Where this sits in CS 300.} Until now, lookup costs $O(n)$
(linear search) or $O(\log n)$ (binary search, balanced trees). Chapter 5
introduces \textbf{hashing}, the technique that takes lookup \emph{on
average} down to $O(1)$. The catch is everywhere: bad hash functions,
full tables, adversarial inputs, and resizing. Mastering hash tables is
mastering the tradeoffs that make $O(1)$ possible.

\textbf{What you add to your toolkit here:}
\begin{itemize}
    \item \textbf{Hash function} as a deterministic $key \to index$ map,
          with the requirements (fast, well-distributed, deterministic).
    \item \textbf{Collision resolution}: chaining (buckets are linked
          lists) vs.\ open addressing (linear / quadratic probing,
          double hashing).
    \item \textbf{Load factor} $\alpha = n / \text{capacity}$: the single
          number that predicts whether your table is fast or sad.
    \item \textbf{Resizing} (rehashing) as the mechanism that keeps
          amortized $O(1)$ honest.
    \item \textbf{Direct hashing} -- the degenerate case where keys
          \emph{are} indices. Works when the key space is dense and
          small.
    \item \textbf{Cryptographic hashing} -- same mathematical object,
          entirely different requirements (collision resistance,
          preimage resistance, slowness). Don't confuse the two.
\end{itemize}

\textbf{Mastery checklist} -- you should be able to, without references:
\begin{enumerate}
    \item Define a hash function and state the three desiderata (fast,
          deterministic, uniform).
    \item Given a key stream, draw the resulting table under chaining
          AND under linear probing, starting from capacity 7.
    \item State the expected cost of insert/find/remove under chaining
          and under open addressing, and say what happens as $\alpha$
          approaches 1.
    \item Compute the mid-square or multiplication-method hash of a
          concrete key by hand.
    \item Explain \emph{primary clustering} (linear probing) and
          \emph{secondary clustering} (quadratic probing) and how double
          hashing escapes both.
    \item Describe the rehash procedure: when it triggers, what the new
          capacity is, and why you can't just copy the old table -- you
          must rehash every element.
    \item Say why \texttt{std::unordered\_map} uses chaining (not open
          addressing) and what that buys vs.\ costs.
    \item Distinguish cryptographic and data-structure hashing, and name
          one attack (e.g., hash flooding) that bridges them.
\end{enumerate}

\textbf{Looking ahead.} Hash tables will show up in graph algorithms as
``visited'' sets (opt.\ ch.~10), in DP as memoization tables, and as the
default \emph{map} container in every production C++ codebase
(\texttt{std::unordered\_map}). Chapter 6 (BSTs) gives you the
\emph{ordered} alternative -- slower on average, but with guarantees
hash tables lack.
\end{ideabox}

\section{5.1 Hash Tables: The Big Idea}

Every data structure you've seen so far has been one of two shapes:
a \emph{list} (search costs $O(n)$) or a \emph{sorted array} (search
costs $O(\log n)$). A hash table does something radically different.
Instead of searching for a key, it \emph{computes} where the key
should live. Search, insert, and remove become $O(1)$ on
average --- constant time, regardless of how many items are
stored.

\begin{defnbox}[Hash table]
A data structure that stores items in an array and uses a
\emph{hash function} to compute, from each item's \emph{key}, the
array index (``bucket'') where the item should be stored. Lookup,
insertion, and removal are therefore one arithmetic computation
plus one array access --- $O(1)$ on average.
\end{defnbox}

\begin{notebox}[Hash tables implement the Set ADT]
Recall the Sequence-vs-Set framing from \textsection 4.1: every
container belongs to one of two interface families. \textbf{Sequence}
keeps items at caller-chosen positions ($i = 0, 1, \ldots, n-1$);
\textbf{Set} keeps items at positions chosen by the data structure
itself, looked up by an \emph{intrinsic key}. ch.~4 built Sequence
implementations (lists, stacks, queues, deques). \emph{This chapter
and ch.~6 build Set implementations.}

A hash table is a Set: the operations are \texttt{find(k)},
\texttt{insert(x)}, \texttt{delete(k)} --- never ``insert at position
$i$'' or ``the third item.'' In C++:
\begin{itemize}[leftmargin=*]
  \item \texttt{std::unordered\_set<K>} --- the bare Set: ``is key
        $k$ present?''
  \item \texttt{std::unordered\_map<K,V>} --- a \emph{dictionary}
        (Set without the ordered ops, in OCW terminology): ``what
        value is associated with key $k$?''
\end{itemize}
ch.~5 covers the chained and open-addressed implementations of the
Set interface. \textsection 5.8 covers the special case where keys
\emph{are} small integers and the hash function is the identity ---
the trivial Set implementation. ch.~6 covers BSTs, the alternative
Set implementation that gives up O(1) lookup in exchange for
ordered operations (\texttt{find\_min}, \texttt{find\_next(k)}).
\end{notebox}

\subsection*{The core trick}

Suppose you have an array of size 10 and you want to store the
key 25. Where does it go?

\begin{quote}
\texttt{index = 25 \% 10 = 5}
\end{quote}

Store it at \texttt{array[5]}. To look it up later: compute the
same index, read that slot. No search. No comparison loop. The
work of ``finding'' has been replaced by a calculation that
locates the element directly.

\begin{examplebox}[Simplest possible hash table]
\begin{lstlisting}[language=C++,basicstyle=\ttfamily\small,frame=single]
constexpr size_t N = 10;
int table[N];

void insert(int key) { table[key % N] = key; }
bool contains(int key) { return table[key % N] == key; }
\end{lstlisting}
\end{examplebox}

Three lines, $O(1)$ in both operations. This is the entire idea.
Everything else in this chapter --- hash functions, collisions,
chaining, probing, load factor --- is about making this trick
\emph{actually work} in practice when keys don't divide up so
neatly.

\subsection*{Why not just use a giant array?}

If keys are small integers, you already can. An array of size
$2^{32}$ indexed directly by \texttt{int} keys works --- but
that's 16 GB of memory to hold possibly a handful of values.
Hash tables compress a large \emph{potential} key space into a
small array, trading perfect placement for practical memory use.

\begin{center}
\begin{tabular}{|l|l|l|}
\hline
\textbf{Approach} & \textbf{Lookup} & \textbf{Space} \\
\hline
Direct array (key = index)       & $O(1)$ & $O(K)$ where K = key range \\
Sorted array + binary search     & $O(\log n)$ & $O(n)$ \\
Linked list                      & $O(n)$ & $O(n)$ \\
Hash table                       & $O(1)^*$ & $O(n)$ \\
\hline
\end{tabular}
\end{center}

\noindent $^*$ average case; worst case can degrade, see later sections.

The hash table gets you array-speed lookup at list-space cost.
That combination is so valuable that hash tables are the backbone
of almost every real-world fast key-value store: JavaScript
objects, Python dicts, Java \texttt{HashMap}, Redis, DNS caches,
database indexes (sometimes), compiler symbol tables, the list
goes on.

\subsection*{Vocabulary}

\begin{center}
\begin{tabular}{|l|l|}
\hline
\textbf{Term} & \textbf{Meaning} \\
\hline
Key            & The value used to look up an item \\
Value          & The data associated with the key (often absent --- just a set) \\
Bucket         & A slot in the underlying array \\
Hash function  & \texttt{key} $\to$ \texttt{bucket index} \\
Collision      & Two keys that hash to the same bucket \\
Load factor    & $n$ items / $m$ buckets --- a density measure \\
\hline
\end{tabular}
\end{center}

\begin{ideabox}[Collisions are the entire problem]
The naive hash table above breaks the instant two keys hash to
the same bucket. With 10 buckets and key \texttt{\% 10}, keys 25
and 75 both land at index 5 --- the second insert either
overwrites the first or needs a plan B. The rest of the chapter
is ``plan B'': different strategies for resolving collisions. The
big split is \emph{chaining} (each bucket holds a list of items)
vs. \emph{open addressing} (find another empty bucket in the
array). Each has different cost/space tradeoffs.
\end{ideabox}

\begin{warnbox}[$O(1)$ is the \emph{average} case]
A hash table's $O(1)$ lookup is the average case under a good
hash function and reasonable load. Worst case, every key hashes
to the same bucket and lookup degrades to $O(n)$ --- the same as
a linked list. Two defenses: pick a good hash function
(distribution matters), and keep the load factor low (resize the
array when it fills up). Sections 5.5--5.8 make this concrete.
\end{warnbox}

\begin{notebox}[Hash tables vs. STL]
C++ gives you two hash-based containers in \texttt{<unordered\_map>}
and \texttt{<unordered\_set>}. Both are chaining hash tables. The
older \texttt{std::map} / \texttt{std::set} are \emph{balanced
BSTs} (red-black trees, chapter 6), which is why their operations
are $O(\log n)$ instead of $O(1)$. ``Ordered'' in the name is the
giveaway: iterating a \texttt{map} gives sorted order; iterating
an \texttt{unordered\_map} gives whatever order the hash function
happened to produce. If you don't need ordering, unordered is
usually faster.
\end{notebox}

\begin{ideabox}[Chapter 5 roadmap]
This chapter is structured as: hash table idea $\to$ hash
functions $\to$ collision resolution (two strategies) $\to$
cost analysis and load factor. By the end, you'll be able to
implement one from scratch and reason about its performance.
The key insight to hold onto: \emph{a hash table is an array
accessed through a function, not through a search loop.}
Everything else is engineering around that idea.
\end{ideabox}

\section{5.2 Chaining}

\emph{Chaining} is the first collision-resolution strategy. The
idea is almost embarrassingly simple: if two keys want the same
bucket, don't pick a winner --- let both live there, in a list.

\begin{defnbox}[Chaining]
A collision-resolution strategy in which each bucket stores a
\emph{list} of items rather than a single item. On insertion,
the new item is appended to the bucket's list; on lookup, the
bucket's list is searched for the matching key.
\end{defnbox}

\subsection*{The mental model}

Picture the hash table as an array of \emph{list heads} instead
of an array of values. Each slot is either \texttt{nullptr}
(empty) or points to a chain of nodes that all hashed to that
index.

\begin{center}
\begin{tabular}{|c|l|}
\hline
\textbf{Bucket} & \textbf{Chain} \\
\hline
0 & \texttt{nullptr} \\
1 & 21 $\to$ 31 $\to$ 51 \\
2 & \texttt{nullptr} \\
3 & 13 \\
4 & 24 $\to$ 74 \\
\hline
\end{tabular}
\end{center}

Every chain only contains keys that hashed to that index, so the
search step only compares against ``nearby'' items --- ideally
just one or two, not the entire table.

\subsection*{Algorithms}

Each of the three operations becomes: (1) hash the key to find
the bucket, (2) do the list-level operation on that bucket's
chain.

\begin{examplebox}[Chained hash table operations]
\begin{lstlisting}[language=C++,basicstyle=\ttfamily\small,frame=single]
template <typename T>
class ChainedHashTable {
    struct Node { T data; Node* next; };
    std::vector<Node*> table;
    size_t n_ = 0;

    size_t hash(const Key& k) const {
        return std::hash<Key>{}(k) % table.size();
    }

public:
    explicit ChainedHashTable(size_t m) : table(m, nullptr) {}

    bool contains(const Key& k) const {
        for (Node* p = table[hash(k)]; p; p = p->next)
            if (p->data.key == k) return true;
        return false;
    }

    void insert(const T& item) {
        size_t b = hash(item.key);
        for (Node* p = table[b]; p; p = p->next)
            if (p->data.key == item.key) { p->data = item; return; }
        table[b] = new Node{item, table[b]};  // prepend
        ++n_;
    }

    bool remove(const Key& k) {
        size_t b = hash(k);
        Node** cur = &table[b];
        while (*cur) {
            if ((*cur)->data.key == k) {
                Node* dead = *cur;
                *cur = dead->next;
                delete dead;
                --n_;
                return true;
            }
            cur = &(*cur)->next;
        }
        return false;
    }
};
\end{lstlisting}
\end{examplebox}

Two implementation notes: \texttt{insert} checks for duplicates
first because most hash-table APIs are ``unique-key'' (insert acts
as upsert), and \texttt{remove} uses the \texttt{Node**} trick
from chapter 4 so the same code handles head and interior removal
without a special case.

\subsection*{Cost}

\begin{center}
\begin{tabular}{|l|c|c|l|}
\hline
\textbf{Operation} & \textbf{Average} & \textbf{Worst} & \textbf{Notes} \\
\hline
\texttt{insert}   & $O(1 + \alpha)$ & $O(n)$ & Chain walk to check duplicate \\
\texttt{contains} & $O(1 + \alpha)$ & $O(n)$ & Chain walk to compare keys \\
\texttt{remove}   & $O(1 + \alpha)$ & $O(n)$ & Chain walk to locate node \\
\hline
\end{tabular}
\end{center}

\noindent $\alpha = n/m$ = \emph{load factor} = average chain
length. When $\alpha$ is a small constant (say, $\leq 1$), every
operation is effectively $O(1)$. When $\alpha$ grows, operations
slow linearly. Section 5.6 discusses how to keep $\alpha$ bounded
by resizing.

\begin{ideabox}[The key insight]
With a good hash function and a reasonably-sized table, chains
stay short --- often length 0, 1, or 2. So even though the
\emph{worst} case is $O(n)$, the \emph{average} case is a
handful of comparisons. The $O(1)$ claim for hash tables is
really ``$O(1 + \alpha)$ where $\alpha$ is kept small by good
engineering.''
\end{ideabox}

\subsection*{Pros and cons}

\begin{center}
\begin{tabular}{|l|p{9cm}|}
\hline
\textbf{Pro} & Never fills up --- you can keep inserting past the
table size, chains just grow longer. \\
\textbf{Pro} & Removal is easy --- just unsplice a list node. \\
\textbf{Pro} & High load factors tolerable (still works at $\alpha > 1$). \\
\textbf{Con} & Extra memory per item (one \texttt{next} pointer per node). \\
\textbf{Con} & Cache-unfriendly --- chain walks follow pointers all
over the heap, same as plain linked lists. \\
\textbf{Con} & Extra allocation on every insert (one \texttt{new Node}). \\
\hline
\end{tabular}
\end{center}

\begin{warnbox}[A degenerate hash function degrades to a list]
If every key hashes to the same bucket, chaining leaves you with
a single linked list of length $n$ --- and every operation is
$O(n)$. This is not hypothetical: some attackers exploit
predictable hash functions to deliberately trigger this
(``hash-collision DoS''). Modern libraries use randomized hash
seeds to make collision attacks impractical. If you roll your
own hash function, be aware that ``obvious'' functions (like
\texttt{key \% m} with adversarial keys) can produce this
worst-case behavior.
\end{warnbox}

\begin{notebox}[Why $O(1+\alpha)$? The simple uniform hashing assumption]
The cost claim ``find takes $O(1+\alpha)$ on average'' isn't a
slogan; it's a theorem (CLRS Ch 11.2, Theorems 11.1--11.2). The
proof rests on one assumption:

\textbf{Simple uniform hashing.} Any given key is equally likely
to hash to any of the $m$ slots, independently of where any other
key has hashed.

Under this assumption, the expected length of bucket $T[h(k)]$ is
$E[n_{h(k)}] = n/m = \alpha$. So an unsuccessful search examines
$\alpha$ elements on average, plus $O(1)$ to compute the hash and
index --- total $\Theta(1 + \alpha)$. A symmetric argument (with
slightly more bookkeeping for the ``which list is searched''
probability) gives the same $\Theta(1 + \alpha)$ bound for
successful search.

Two consequences worth internalizing:
\begin{itemize}[leftmargin=*]
  \item If you keep $\alpha$ bounded by a constant (which is what
        resizing in \textsection 5.6 does), the $1 + \alpha$ term
        is $O(1)$. \emph{This is the entire reason hash tables
        average $O(1)$.}
  \item The open-addressing variants in \textsection 5.3--\textsection 5.5
        all use the same simple-uniform-hashing assumption. Their
        per-strategy cost formulas (linear's
        $\frac{1}{2}(1 + 1/(1-\alpha))$, double-hashing's tighter
        bounds) come from analyzing different probe sequences
        \emph{under that same assumption}.
\end{itemize}
The assumption is strong --- real hash functions don't perfectly
satisfy it on all inputs. \textsection 5.7 shows what hash
functions try to do to approximate it; \textsection 5.7 (universal
hashing subsection) shows how to make a probabilistic guarantee
without the assumption at all.
\end{notebox}

\begin{notebox}[What the STL does]
\texttt{std::unordered\_map} and \texttt{std::unordered\_set} are
\emph{required} by the standard to use chaining (separate
chaining, in formal terminology). This is a correctness
requirement, not just an implementation choice: the standard
guarantees that inserting doesn't invalidate iterators to other
elements, which essentially forces per-bucket lists. The
alternative, open addressing (next section), moves elements
around on insertion and would break iterator stability.
\end{notebox}

\begin{ideabox}[The contrast with open addressing]
Chaining solves collisions by \emph{going outside the array} ---
add extra memory in the form of chains. The next section covers
the opposite approach: \emph{stay inside the array} and find a
different bucket for the collider. Both get $O(1)$ average, with
different space/cache tradeoffs. Know both; the choice depends
on context.
\end{ideabox}

\section{5.3 Linear Probing}

\emph{Linear probing} is the first flavor of \emph{open
addressing} --- the ``stay inside the array'' approach to
collisions. When the computed bucket is occupied, walk forward
through the array one slot at a time until you find an empty
one. That's it. Simple to describe, surprisingly subtle to get
right.

\begin{defnbox}[Linear probing]
A collision-resolution strategy that stores all items directly in
the hash table array. On collision, it checks bucket $h+1$, then
$h+2$, then $h+3$, \ldots (wrapping around the end), until an
empty bucket is found. The sequence $h, h+1, h+2, \ldots$ is
called the \emph{probe sequence}.
\end{defnbox}

\subsection*{Why it's interesting}

Chaining (5.2) allocates a new node per item. Linear probing does
\emph{not} --- every item lives in the fixed array itself.

\begin{center}
\begin{tabular}{|l|l|l|}
\hline
\textbf{Aspect} & \textbf{Chaining} & \textbf{Linear probing} \\
\hline
Extra memory per item & 1 pointer  & 0 \\
Allocations on insert & 1 \texttt{new Node} & 0 \\
Cache behavior        & Poor (heap jumps) & Excellent (contiguous) \\
Can overflow?         & Never     & Yes (array full $\Rightarrow$ insert fails) \\
Removal               & Unsplice list node & Tombstone (see below) \\
\hline
\end{tabular}
\end{center}

For small-to-medium tables, linear probing typically \emph{wins}
on real hardware: no allocation, no pointer chasing, and the CPU
prefetches the next probe slot because it's literally the next
cache line. This is why many high-performance hash tables (Google's
\texttt{dense\_hash\_map}, Robin Hood hashing, Swiss tables) are
built on open addressing rather than chaining.

\subsection*{Insert}

\begin{examplebox}[Linear-probe insert]
\begin{lstlisting}[language=C++,basicstyle=\ttfamily\small,frame=single]
bool insert(int key) {
    size_t start = hash(key) % N;
    for (size_t i = 0; i < N; ++i) {
        size_t b = (start + i) % N;
        if (table[b].state != OCCUPIED) {
            table[b] = { key, OCCUPIED };
            return true;
        }
        if (table[b].key == key) return false; // duplicate
    }
    return false;  // table full
}
\end{lstlisting}
\end{examplebox}

The \texttt{\% N} wrap makes the probe sequence circular --- if
you hit the end of the array, you keep going from index 0.

\subsection*{The tombstone problem}

Removal on an open-addressed table is \emph{not} just ``mark the
bucket empty.'' Consider:

\begin{enumerate}[leftmargin=*]
    \item Insert 25 (bucket 5).
    \item Insert 35 (bucket 5 occupied, probe to 6).
    \item Remove 25 --- mark bucket 5 empty.
    \item Search for 35 --- check bucket 5, find it empty, stop
          searching. Returns ``not found.'' \textbf{But 35 is
          there}, in bucket 6.
\end{enumerate}

To fix this, open-addressed tables distinguish \emph{two} kinds
of empty:

\begin{center}
\begin{tabular}{|l|l|}
\hline
\textbf{State} & \textbf{Meaning} \\
\hline
\texttt{EMPTY\_SINCE\_START} & Never held an item; search can stop here \\
\texttt{OCCUPIED}            & Currently holds an item \\
\texttt{EMPTY\_AFTER\_REMOVAL} (tombstone) & Was occupied; search continues past \\
\hline
\end{tabular}
\end{center}

The tombstone preserves the ``I was here'' information that the
probe sequence needs. Insertions treat tombstones as empty (can
reuse the slot). Searches treat tombstones as occupied (keep
probing past).

\begin{warnbox}[Tombstones never clean themselves up]
A long-lived table with many insertions and removals can fill up
with tombstones, slowing searches to a crawl. The fix is to
\emph{rehash}: allocate a fresh table, re-insert every live item,
throw the old table away. The tombstones are gone, probes are
short again. Sections 5.6--5.8 discuss when and how to trigger
this.
\end{warnbox}

\subsection*{Search}

\begin{examplebox}[Linear-probe search]
\begin{lstlisting}[language=C++,basicstyle=\ttfamily\small,frame=single]
int* find(int key) {
    size_t start = hash(key) % N;
    for (size_t i = 0; i < N; ++i) {
        size_t b = (start + i) % N;
        if (table[b].state == EMPTY_SINCE_START) return nullptr;
        if (table[b].state == OCCUPIED && table[b].key == key)
            return &table[b].key;
        // tombstone: keep looking
    }
    return nullptr;
}
\end{lstlisting}
\end{examplebox}

The three termination conditions matter: (1) found a match,
return it; (2) hit a never-occupied slot, definitely not present;
(3) probed every bucket, not present.

\subsection*{Cost and load factor}

Linear probing's cost depends dramatically on how full the table
is. Let $\alpha = n/m$ (load factor).

\begin{center}
\begin{tabular}{|l|c|l|}
\hline
\textbf{Load factor $\alpha$} & \textbf{Avg probes (successful)} & \textbf{Feel} \\
\hline
0.5  & $\approx 1.5$  & Fast \\
0.75 & $\approx 2.5$  & Still OK \\
0.9  & $\approx 5.5$  & Getting sluggish \\
0.95 & $\approx 10.5$ & Noticeable \\
0.99 & $\approx 50$   & Disaster \\
\hline
\end{tabular}
\end{center}

These numbers aren't linear in $\alpha$ --- they grow like
$\frac{1}{1 - \alpha}$ (Knuth's analysis). That's why open-
addressed tables typically enforce a much lower max load factor
than chained ones: 0.5--0.75 is common for linear probing versus
1.0+ for chaining.

\begin{ideabox}[Primary clustering]
Linear probing has a pathological behavior called
\emph{primary clustering}: once a run of occupied slots forms,
any key that hashes \emph{into} that run will extend it, making
it more likely to catch future keys too. Clusters grow like
snowballs. This is why quadratic probing (5.4) and double hashing
(5.5) exist --- they break the cluster-growth pattern by making
probe sequences depend on the key, not just the starting bucket.
\end{ideabox}

\begin{notebox}[Why linear probing is still worth using]
Despite primary clustering, linear probing often beats ``smarter''
probe sequences in practice because of cache behavior. Every
probe is the next cache line --- no pointer chasing, no random
access. Quadratic probing and double hashing have
theoretically-shorter probe sequences but jump around the array,
busting the cache. At load factors under 0.75, linear probing's
wall-clock speed is hard to beat.
\end{notebox}

\begin{ideabox}[Key takeaways]
Open addressing trades ``cheap removal'' for ``no allocation''
and ``great cache behavior.'' The tombstone mechanism is the
essential correctness fix. Load factor discipline is the
essential performance fix --- probe counts explode near $\alpha =
1$. The next sections show variants of probing that keep the
same flavor but sidestep primary clustering.
\end{ideabox}

\section{5.4 Quadratic Probing}

Linear probing's flaw, remember, is \emph{primary clustering}:
once a run of occupied buckets forms, any key that hashes into
that run lengthens it, snowballing the problem. \emph{Quadratic
probing} sidesteps clustering by using a probe sequence whose
gaps grow with each step.

\begin{defnbox}[Quadratic probing]
An open-addressing collision-resolution strategy where the probe
sequence for a key $k$ is
$$\text{bucket}_i = (h(k) + c_1 i + c_2 i^2) \bmod m$$
for $i = 0, 1, 2, \ldots$ and chosen constants $c_1, c_2$. The
quadratic term makes successive probes spread out instead of
walking contiguously.
\end{defnbox}

\subsection*{The probe pattern}

With $c_1 = 0, c_2 = 1$ (the simplest common choice), the probe
sequence from a starting bucket $h$ visits:

\begin{center}
\begin{tabular}{|c|c|c|}
\hline
$i$ & Offset $i^2$ & Bucket \\
\hline
0 & 0  & $h$ \\
1 & 1  & $h+1$ \\
2 & 4  & $h+4$ \\
3 & 9  & $h+9$ \\
4 & 16 & $h+16$ \\
\hline
\end{tabular}
\end{center}

The gaps grow: 1, 3, 5, 7, 9, \ldots The probe sequence jumps
past nearby occupied slots quickly, so even if there's a cluster
starting near $h$, it gets skipped over rather than extended.

\begin{examplebox}[Quadratic-probe insert]
\begin{lstlisting}[language=C++,basicstyle=\ttfamily\small,frame=single]
bool insert(int key) {
    size_t h = hash(key);
    for (size_t i = 0; i < N; ++i) {
        size_t b = (h + c1 * i + c2 * i * i) % N;
        if (table[b].state != OCCUPIED) {
            table[b] = { key, OCCUPIED };
            return true;
        }
        if (table[b].key == key) return false; // duplicate
    }
    return false;  // full or unreachable
}
\end{lstlisting}
\end{examplebox}

Everything else (tombstones, search, remove) carries over from
linear probing. The only change is the probe-step formula.

\subsection*{Primary clustering: solved. Secondary clustering: new problem}

Quadratic probing eliminates primary clustering (adjacent keys no
longer share probe paths). But two keys that hash to the same
\emph{starting} bucket will still follow the exact same probe
sequence --- that's \emph{secondary clustering}. Less severe than
primary, but not gone.

\begin{center}
\begin{tabular}{|l|c|c|}
\hline
\textbf{Strategy} & \textbf{Primary clustering} & \textbf{Secondary clustering} \\
\hline
Linear probing    & Yes  & Yes \\
Quadratic probing & No   & Yes \\
Double hashing (5.5) & No & No \\
\hline
\end{tabular}
\end{center}

\begin{warnbox}[Quadratic probing may not visit every bucket]
With a poorly-chosen $N$ and constants $c_1, c_2$, the probe
sequence can cycle through a small set of buckets forever and
miss the rest of the table --- insert fails even though empty
buckets exist. There's a theorem: if $N$ is prime and the load
factor is less than 0.5, the probe sequence \emph{is} guaranteed
to visit $N/2$ distinct buckets. So: pick prime $N$, keep
$\alpha \leq 0.5$. If you need higher load factors, use linear
probing or double hashing instead.
\end{warnbox}

\subsection*{Choosing constants}

$c_1 = 0, c_2 = 1$ gives you $h + i^2$, which has the guarantee
above. $c_1 = 1, c_2 = 1$ gives $h + i + i^2$, which is the form
used in the textbook pseudocode; it still works but the theoretical
guarantees depend on specific $N$ values. When in doubt, prefer
the simpler $h + i^2$ version with prime table sizes.

\begin{examplebox}[Triangular-number probing]
A well-known trick uses $c_1 = c_2 = 1/2$, giving probes at
offsets $\frac{i(i+1)}{2}$ (the triangular numbers: 0, 1, 3, 6,
10, 15, \ldots). This form visits \emph{every} bucket when $N$
is a power of 2 --- a major convenience because hash tables often
use power-of-2 sizes for fast modulo via bit masking. This is the
form used in Java's older hash tables and many game-engine hash
maps.
\end{examplebox}

\subsection*{Cost}

\begin{center}
\begin{tabular}{|l|c|l|}
\hline
\textbf{Load factor $\alpha$} & \textbf{Avg probes (successful)} & \textbf{Notes} \\
\hline
0.5  & $\approx 1.4$  & Slightly better than linear at 0.5 \\
0.75 & $\approx 2.0$  & Noticeably better \\
0.9  & $\approx 3.2$  & Much better than linear's 5.5 \\
\hline
\end{tabular}
\end{center}

At high load factors, quadratic probing substantially outperforms
linear probing. At low load factors, the difference is minor
and linear's cache-friendliness can win on wall-clock time.

\begin{ideabox}[Quadratic vs linear: the real tradeoff]
Linear probing: simple, cache-friendly, suffers at high load.
Quadratic probing: spreads clusters, works better at moderate
load, slightly less cache-friendly (the jumps in the probe
sequence skip cache lines). In modern hash table implementations
that push load factors close to 1, the cache penalty is worth
paying to avoid primary clustering. For textbook problems and
moderate loads, either is fine.
\end{ideabox}

\begin{notebox}[Nobody in production uses ``pure'' quadratic probing]
Real-world implementations combine ideas: Google's Abseil Swiss
tables use linear probing with SSE-accelerated parallel
comparison; Rust's \texttt{HashMap} uses Robin Hood hashing (a
linear-probing variant that equalizes probe lengths); CPython's
\texttt{dict} uses a perturbation scheme that's neither pure
linear nor pure quadratic. The textbook probes are pedagogy; the
real insight is ``open addressing plus anything that breaks
clustering.''
\end{notebox}

\section{5.5 Double Hashing}

Linear and quadratic probing both suffer \emph{secondary
clustering}: any two keys that hash to the same starting bucket
follow the same probe sequence. \emph{Double hashing} fixes this
by making the probe \emph{step size} also depend on the key ---
so two keys starting at the same bucket diverge immediately.

\begin{defnbox}[Double hashing]
An open-addressing strategy that uses two hash functions. The
probe sequence for a key $k$ is
$$\text{bucket}_i = (h_1(k) + i \cdot h_2(k)) \bmod m$$
for $i = 0, 1, 2, \ldots$. $h_1$ determines the starting bucket;
$h_2$ determines the step size between probes.
\end{defnbox}

\subsection*{The idea, geometrically}

Think of each key as having its own ``rhythm'' for walking
through the table. Linear probing: everyone steps by 1.
Quadratic: everyone's step pattern is $1, 4, 9, \ldots$.
\emph{Double hashing: each key's step depends on the key itself.}

Two keys can collide on their starting bucket and \emph{still}
separate instantly because their $h_2$ values differ. That's the
whole trick.

\subsection*{Example}

Let $m = 11$, $h_1(k) = k \bmod 11$, $h_2(k) = 5 - (k \bmod 5)$.

\begin{center}
\begin{tabular}{|c|c|c|l|}
\hline
\textbf{Key $k$} & $h_1(k)$ & $h_2(k)$ & \textbf{Probe sequence} \\
\hline
55  & 0   & 5 & 0, 5, 10, 4, 9, 3, \ldots \\
66  & 0   & 4 & 0, 4, 8, 1, 5, \ldots \\
\hline
\end{tabular}
\end{center}

Both keys start at bucket 0, but after the first collision they
walk in entirely different directions. That's why double hashing
approximates ``uniform hashing'' --- the theoretical ideal where
every key's probe sequence is a random permutation of buckets.

\subsection*{Implementation}

\begin{examplebox}[Double-hashing insert]
\begin{lstlisting}[language=C++,basicstyle=\ttfamily\small,frame=single]
bool insert(int key) {
    size_t h1 = hash1(key) % N;
    size_t h2 = hash2(key);   // must be nonzero, coprime to N
    for (size_t i = 0; i < N; ++i) {
        size_t b = (h1 + i * h2) % N;
        if (table[b].state != OCCUPIED) {
            table[b] = { key, OCCUPIED };
            return true;
        }
        if (table[b].key == key) return false; // duplicate
    }
    return false;  // table full
}
\end{lstlisting}
\end{examplebox}

Search and remove use the same probe formula with the usual
tombstone rules. The only new concern is \emph{choosing $h_2$
correctly}.

\subsection*{The $h_2$ requirement: never zero, always coprime to $m$}

If $h_2(k) = 0$, the probe sequence stays stuck at $h_1(k)$
forever --- infinite loop. Always enforce $h_2(k) \neq 0$.

If $h_2(k)$ shares a factor with $m$ (table size), the probe
sequence visits only a subset of buckets --- you can fail to
insert even with empty buckets available. The standard fix:
\textbf{make $m$ a prime and design $h_2$ to always return a
value in $[1, m-1]$}. A common construction:

$$h_2(k) = p - (k \bmod p), \text{ where } p < m \text{ is prime}$$

This guarantees $1 \leq h_2(k) \leq p < m$, and since $m$ is
prime, $\gcd(h_2(k), m) = 1$ always.

\begin{warnbox}[Powers of 2 are incompatible with double hashing]
If $m$ is a power of 2, the only values coprime to $m$ are the
odd ones. So $h_2$ has to always return an odd number --- often
enforced by $h_2(k) = 2 \cdot g(k) + 1$. This works but limits
$h_2$'s range. Most double-hashing implementations use prime
table sizes for this reason, even though modulo-by-prime is
slower than modulo-by-power-of-2.
\end{warnbox}

\subsection*{Cost}

\begin{center}
\begin{tabular}{|l|c|l|}
\hline
\textbf{Load factor $\alpha$} & \textbf{Avg probes (successful)} & \textbf{Notes} \\
\hline
0.5  & $\approx 1.4$   & Matches quadratic \\
0.75 & $\approx 1.8$   & Better than quadratic \\
0.9  & $\approx 2.6$   & Significantly better \\
0.95 & $\approx 3.2$   & Holds up where others collapse \\
\hline
\end{tabular}
\end{center}

The cost formula is $\frac{1}{\alpha} \ln \frac{1}{1-\alpha}$ --
essentially the theoretical best achievable by any open-
addressing scheme. Double hashing's probe distribution is
statistically close to uniform random, which is what the analysis
assumes.

\begin{center}
\begin{tabular}{|l|c|c|c|}
\hline
\textbf{Strategy} & \textbf{Primary clusters?} & \textbf{Secondary clusters?} & \textbf{Cache-friendly?} \\
\hline
Chaining          & ---  & ---  & No \\
Linear probing    & Yes  & Yes  & Yes \\
Quadratic probing & No   & Yes  & Somewhat \\
Double hashing    & No   & No   & No \\
\hline
\end{tabular}
\end{center}

\begin{ideabox}[Why double hashing isn't always used]
Double hashing has the best theoretical probe counts, yet most
modern hash table implementations use linear probing. The reason
is cache. Double hashing's jumps across the table are random
accesses from the CPU's point of view; linear probing's contiguous
walk feeds the cache. At load factors below 0.8, linear probing's
cache advantage often beats double hashing's shorter probe
count in wall-clock time. Above 0.8, double hashing wins. The
lesson: big-O is not the whole cost story.
\end{ideabox}

\begin{notebox}[Two hash functions from one]
A common trick when you only have one hash function: split its
output. For a 64-bit hash $H(k)$, use the low 32 bits as $h_1$
and the high 32 bits as $h_2$. This gives you two
statistically-independent-enough hash values from one computation.
Java's \texttt{HashMap} uses a variant of this idea.
\end{notebox}

\begin{ideabox}[The pattern across 5.2--5.5]
Four collision strategies, one theme: \emph{spread the keys out
so probe paths are short}. Chaining spreads into separate lists.
Linear probing walks, and clusters. Quadratic widens the walk.
Double hashing randomizes the walk. Each fixes a specific flaw
of the previous. Real-world hashing takes the best pieces of
each: open-addressed layout (for cache), randomized probing (for
spread), low load factors (for speed), and periodic rehashing
(for cleanup). The remaining sections tie these together.
\end{ideabox}

\section{5.6 Hash Table Resizing}

Every cost analysis so far has assumed the load factor stays
bounded. But the user keeps inserting, and $n$ keeps growing.
Without intervention, $\alpha = n/m$ climbs past 0.5, 0.75, 0.9,
and eventually the table fills up. \emph{Resizing} is how hash
tables bound their load factor at runtime.

\begin{defnbox}[Resize (rehash)]
The operation of allocating a larger array, reinserting every
existing item into it using a fresh hash-to-bucket mapping, and
discarding the old array. Triggered when the load factor crosses
a threshold. Cost: $O(n)$ --- every item must be rehashed.
\end{defnbox}

\subsection*{The algorithm}

\begin{enumerate}[leftmargin=*]
    \item Decide a new size $m'$. Commonly $m' \approx 2m$, rounded
          up to the next prime.
    \item Allocate a fresh array of size $m'$, all slots empty.
    \item For each non-empty bucket in the old table, re-compute
          $h(key) \bmod m'$ (note: \emph{not} $\bmod m$) and
          insert into the new table.
    \item Free the old array. Swap the pointer.
\end{enumerate}

\begin{examplebox}[Resize]
\begin{lstlisting}[language=C++,basicstyle=\ttfamily\small,frame=single]
void resize(size_t newCap) {
    std::vector<Slot> oldTable = std::move(table);
    table.assign(newCap, Slot{});   // all EMPTY_SINCE_START
    n_ = 0;
    for (const Slot& s : oldTable) {
        if (s.state == OCCUPIED) insert(s.key);
    }
}
\end{lstlisting}
\end{examplebox}

The old bucket indices are \emph{meaningless} for the new table.
Every item must be re-hashed from scratch. This is why the
operation is $O(n)$ rather than ``just copy the array.''

\subsection*{When to resize}

Three common triggers:

\begin{center}
\begin{tabular}{|l|l|l|}
\hline
\textbf{Trigger} & \textbf{When} & \textbf{Typical threshold} \\
\hline
Load factor & $\alpha \geq t$             & 0.5--0.75 (open addr); 1.0+ (chaining) \\
Chain length (chaining) & longest chain $\geq k$ & 4--8 \\
Probe count (open addr) & probes $\geq k$   & $N/3$ or similar \\
\hline
\end{tabular}
\end{center}

Load factor is the standard trigger because it's cheap to check
(just compare \texttt{n\_ / m} to the threshold after each
insert). The other triggers are more responsive to pathological
cases --- a few really-bad buckets can justify resizing even if
$\alpha$ looks okay.

\subsection*{Growth factor: why doubling?}

Suppose you grow by adding a fixed amount ($m' = m + c$) each
time. Inserting $n$ items triggers $n/c$ resizes, each $O(n)$ on
average, so total work is $O(n^2 / c) = O(n^2)$. \textbf{Back to
amortized $O(n)$ per insert on average} --- a disaster.

With $m' = 2m$ (doubling), you resize $O(\log n)$ times over the
life of the table. Each resize $i$ rehashes roughly $2^i$ items.
Total rehash work: $1 + 2 + 4 + \ldots + n \approx 2n = O(n)$.
Amortized $O(1)$ per insert.

\begin{ideabox}[Same amortization as \texttt{std::vector}]
This is the same geometric-growth argument that makes
\texttt{std::vector::push\_back} amortized $O(1)$. The rule
generalizes: any data structure that occasionally needs to
rebuild itself should grow by a constant \emph{factor}, not a
constant \emph{amount}. Array lists, hash tables, deques, even
JIT code caches --- all grow geometrically for this reason.
\end{ideabox}

\subsection*{Why prime sizes?}

If table size $m$ has small factors, keys that share those
factors will cluster. Example: $m = 10$ and keys are even. Every
even key hashes to an even bucket; half the table is wasted.

Prime sizes avoid this: $\gcd(\text{key}, m) = 1$ for most keys,
so the modulo spreads them uniformly regardless of key
structure. Doubling to the \emph{next prime} ($m = 11 \to 23 \to
47 \to 97$) is a common implementation choice.

\begin{warnbox}[Resize cost is not invisible]
A single insertion that triggers a resize is $O(n)$ --- the
amortized analysis doesn't help the caller who just hit the
spike. For real-time code (games, audio, trading systems), this
is a problem: one unlucky insert blows the frame budget. Fixes
include: pre-size the table with \texttt{reserve()}, use
incremental rehashing (move a few items per operation over time
instead of all at once --- Redis does this), or accept higher
memory and use a never-resize table.
\end{warnbox}

\subsection*{Shrink?}

Most implementations \emph{never} shrink automatically. Why:

\begin{itemize}[leftmargin=*]
    \item Shrinking and regrowing in a tight loop of inserts and
          removes would thrash.
    \item The memory ``wasted'' by a half-empty table is usually
          cheap compared to the cost of a resize.
    \item If the user really wants to reclaim memory, they can
          call \texttt{rehash(n)} or move into a fresh table
          explicitly.
\end{itemize}

\texttt{std::unordered\_map} has \texttt{rehash(n)} and
\texttt{reserve(n)} but no auto-shrink. \texttt{std::vector} has
\texttt{shrink\_to\_fit()} but it's still a manual request.

\begin{notebox}[Load factor thresholds in the wild]
\texttt{std::unordered\_map} default max: 1.0. Java
\texttt{HashMap}: 0.75. Python \texttt{dict}: 2/3 (0.67). Google
Abseil Swiss tables: $7/8$ (0.875). Rust \texttt{HashMap}: 7/8.
Chaining tables can tolerate $\alpha > 1$ because longer chains
only hurt linearly; open-addressed tables can't, so they resize
earlier. The precise threshold is an engineering knob with
tradeoffs on both sides (too low: wasted memory; too high: slow
probes).
\end{notebox}

\begin{ideabox}[Resizing is what makes ``$O(1)$ average'' work]
Without resizing, a hash table's performance degrades linearly as
you insert more items. The $O(1)$ guarantee only holds because
the implementation keeps $\alpha$ bounded by rebuilding.
Resizing is the \emph{mechanism} that preserves the asymptotic
guarantee --- without it, the analysis breaks down. When you
profile a hash table and see occasional slow operations, it's
almost always a resize in progress.
\end{ideabox}

\section{5.7 Common Hash Functions}

Every hash table so far has assumed a hash function that ``maps
keys to buckets uniformly.'' This section finally makes that
concrete: what does a good hash function \emph{look like}, and
what can go wrong?

\begin{defnbox}[Good hash function]
A function $h: \text{Keys} \to [0, m)$ that (1) is cheap to
compute, (2) distributes keys uniformly across buckets, and (3)
uses \emph{all} the bits of the key so that small differences
between keys produce different hashes. A \emph{perfect} hash
function maps each expected key to a unique bucket with no
collisions --- achievable only when the full key set is known
up front.
\end{defnbox}

\subsection*{Modulo: the simplest, surprisingly subtle}

$h(k) = k \bmod m$ is the textbook default.

\begin{examplebox}[Modulo hash]
\begin{lstlisting}[language=C++,basicstyle=\ttfamily\small,frame=single]
size_t hashModulo(int key, size_t m) {
    return key % m;
}
\end{lstlisting}
\end{examplebox}

When does it fail? When the keys share a factor with $m$. If
keys are always even and $m$ is even, odd buckets are never
used. If $m = 10$ and keys are \texttt{10, 20, 30, 40}, every
one hashes to bucket 0. This is exactly why prime table sizes
are recommended: $\gcd(\text{key}, m)$ is almost always 1, so
the modulo spreads the keys out regardless of structure.

\begin{warnbox}[Powers of 2 and modulo]
$k \bmod 2^n$ reduces to $k$ \texttt{\& ((1 << n) - 1)} --- just
the low $n$ bits. That's fast, but it throws away all the
high-order information. If your keys are clustered in their low
bits (pointers on 64-bit systems often are), collisions explode.
Either use prime $m$ (accept slower modulo) or mix the bits
first (multiply-and-shift, XOR fold).
\end{warnbox}

\subsection*{Mid-square}

Square the key, extract the middle bits. The idea: squaring mixes
every bit of the key into the middle of the result, so the middle
bits depend on the entire key.

\begin{examplebox}[Mid-square hash (binary)]
\begin{lstlisting}[language=C++,basicstyle=\ttfamily\small,frame=single]
size_t hashMidSquare(uint32_t key, size_t m) {
    uint64_t sq   = uint64_t(key) * key;
    size_t   r    = 10;                   // bits to extract
    size_t   drop = (64 - r) / 2;
    uint64_t mid  = (sq >> drop) & ((1ULL << r) - 1);
    return mid % m;
}
\end{lstlisting}
\end{examplebox}

Mid-square is a historical curiosity now --- multiplicative
hashing (below) dominates for integer keys --- but it's a useful
illustration of the general technique: \emph{mix the bits, then
take a slice.}

\subsection*{Multiplicative hashing for strings}

For string keys, you can't just take a single character and
modulo. You need a function that combines every character into a
single hash.

\begin{examplebox}[djb2, Bernstein's classic]
\begin{lstlisting}[language=C++,basicstyle=\ttfamily\small,frame=single]
size_t hashString(const std::string& s) {
    size_t h = 5381;
    for (char c : s) h = h * 33 + static_cast<unsigned char>(c);
    return h;
}
\end{lstlisting}
\end{examplebox}

The idea: iteratively fold each character into a running hash by
multiplying and adding. The multiplier (33) and seed (5381) are
chosen so that small changes in the input produce large changes
in the hash --- an informal \emph{avalanche} property.

Variants:
\begin{itemize}[leftmargin=*]
    \item djb2 XOR: \texttt{h = h * 33 \textasciicircum{} c} --- marginally better distribution
    \item sdbm: \texttt{h = c + (h << 6) + (h << 16) - h} --- good for text
    \item FNV-1a: iterative XOR-then-multiply --- widely used
    \item CityHash, FarmHash, xxHash: modern, SIMD-accelerated, extremely fast
\end{itemize}

\subsection*{What makes a hash function ``good''?}

\begin{center}
\begin{tabular}{|l|p{10cm}|}
\hline
\textbf{Property} & \textbf{Meaning} \\
\hline
Determinism & Same input $\Rightarrow$ same output, always \\
Uniform distribution & Each bucket receives roughly $n/m$ keys \\
Avalanche & Flipping one input bit flips about half the output bits \\
Cheap     & A single arithmetic operation plus modulo, ideally \\
Non-trivial & Doesn't fall into obvious traps (e.g., identity function) \\
\hline
\end{tabular}
\end{center}

``Cryptographic strength'' is usually \emph{not} required. SHA-256
is a good hash function for ensuring data integrity, but
terrible for hash tables (thousands of cycles per call). Hash
table hash functions need speed over security.

\begin{warnbox}[Don't invent your own]
Writing a good hash function is harder than it looks. Subtle
distribution biases only show up under specific workloads and
can be hard to detect. For production code, use
\texttt{std::hash<T>} (C++), Google's CityHash, or xxHash ---
all battle-tested. Hand-rolling ``just use modulo'' is fine for
homework; do not ship it in real software handling adversarial
input.
\end{warnbox}

\subsection*{C++: \texttt{std::hash}}

C++ provides \texttt{std::hash<T>} for the standard types: ints,
pointers, strings, and enums. You can specialize it for your own
types by defining \texttt{std::hash<YourType>} in the
\texttt{std} namespace, or supply a custom hasher as a template
parameter to the unordered container.

\begin{examplebox}[Custom hasher]
\begin{lstlisting}[language=C++,basicstyle=\ttfamily\small,frame=single]
struct Point { int x, y; };

struct PointHash {
    size_t operator()(const Point& p) const noexcept {
        // combine two hashes --- the standard "boost::hash_combine" pattern
        size_t h1 = std::hash<int>{}(p.x);
        size_t h2 = std::hash<int>{}(p.y);
        return h1 ^ (h2 + 0x9e3779b9 + (h1 << 6) + (h1 >> 2));
    }
};

std::unordered_set<Point, PointHash> points;
\end{lstlisting}
\end{examplebox}

The magic number \texttt{0x9e3779b9} is $\frac{2^{32}}{\phi}$
($\phi$ = golden ratio) --- a widely-used bit-mixing constant
that gives good distribution across field combinations. Boost's
\texttt{hash\_combine} uses this; many custom hashers copy the
idiom.

\begin{notebox}[Hash randomization: security vs. reproducibility]
Some languages randomize their hash seeds at startup to defend
against hash-collision DoS attacks (Python 3.3+, Ruby, Rust). A
side effect: iteration order of a hash map is now
\emph{nondeterministic} between program runs. Code that depends
on consistent iteration order breaks subtly. C++'s standard
doesn't mandate randomization, though some implementations offer
it as a build option.
\end{notebox}

\begin{ideabox}[The big picture]
A hash table's speed is only as good as its hash function. Any
of the collision-resolution schemes (chaining, linear, quadratic,
double) will perform terribly if the hash function is poor ---
because they all assume approximately-uniform distribution. When
your hash table is slow, suspect the hash function first, the
load factor second, and the collision strategy last.
\end{ideabox}

\subsection*{Universal hashing: the principled defense against adversarial input}

The hash randomization notebox above is a practical sketch; the
underlying theory is \emph{universal hashing}, due to Carter and
Wegman (1979). It is the standard answer to ``what if my hash
function's worst-case input shows up?'' --- a question every fixed
hash function leaves dangling.

\begin{defnbox}[Universal hash family]
A finite collection $\mathcal{H}$ of hash functions
$h: U \to \{0, 1, \ldots, m-1\}$ is \textbf{universal} if for every
pair of distinct keys $k, l \in U$,
\[
\Pr_{h \in \mathcal{H}}\bigl[h(k) = h(l)\bigr] \;\leq\; \frac{1}{m}.
\]
The probability is over the random choice of $h$, \emph{not} over
the keys.
\end{defnbox}

The motivation is the gap between average-case and worst-case for
\emph{any single} hash function. Pick any fixed $h$ and an
adversary can find $n$ keys that all hash to the same bucket,
giving $\Theta(n)$ chains. Universal hashing escapes this by
randomizing the function itself at the start of execution: the
adversary picks the keys, but the program picks the hash function
\emph{independently and after} the keys are chosen, so worst-case
inputs no longer exist for any particular run.

\begin{examplebox}[The $h_{a,b}$ family from CLRS Ch 11.3.3]
\begin{lstlisting}[language=C++,basicstyle=\ttfamily\small,frame=single]
// Universal hash family h_{a,b}(k) = ((a*k + b) mod p) mod m
// where p > u (universe size) is a fixed prime, and a, b are
// chosen at random at startup with a in [1, p-1], b in [0, p-1].

struct UniversalHash {
    uint64_t a, b;
    uint64_t p;            // prime > universe size
    size_t   m;            // table size

    UniversalHash(size_t table_size, uint64_t prime) : p(prime), m(table_size) {
        std::mt19937_64 rng(std::random_device{}());
        a = 1 + rng() % (p - 1);   // a in [1, p-1], nonzero
        b =     rng() %  p;        // b in [0, p-1]
    }

    size_t operator()(uint64_t k) const {
        return ((a * k + b) % p) % m;
    }
};
\end{lstlisting}
\end{examplebox}

\begin{notebox}[Why universal hashing implies short chains]
The expected chain length at $h(k)$ is now bounded by
$1 + \alpha$ \emph{regardless of the keys}. Sketch (CLRS Theorem
11.3, OCW lec4): for each pair $k, l$ of distinct keys, define
indicator $X_{kl} = \mathbf{1}[h(k) = h(l)]$. By the universality
condition, $E[X_{kl}] = \Pr[h(k)=h(l)] \leq 1/m$. The number of
keys colliding with $k$ is $Y_k = \sum_{l \neq k} X_{kl}$, so by
linearity of expectation,
\[
E[Y_k] \;=\; \sum_{l \neq k} E[X_{kl}] \;\leq\; \sum_{l \neq k} \tfrac{1}{m}
       \;\leq\; \tfrac{n-1}{m} \;\leq\; \alpha.
\]
The expected length of the bucket containing $k$ is $E[Y_k] + 1
\leq 1 + \alpha$. \emph{The same bound as simple uniform hashing,
but now it follows from the algorithm's randomness, not from an
assumption about the input.}
\end{notebox}

\begin{ideabox}[How real systems use this]
Few production hash tables ship a full universal-family
implementation; the constant factors don't justify it for typical
workloads. What real systems \emph{do} ship is the same idea in a
weaker form: a process-startup random seed combined with a
fast non-cryptographic hash. Python's \texttt{hash} (since 3.3),
Rust's \texttt{RandomState}, Java's salted \texttt{String.hashCode}
on recent JVMs --- all of them randomize at process start so an
attacker cannot precompute colliding inputs. The formal
universal-family construction is the textbook answer; salted-hash
randomization is the production answer. The defense is the same
in spirit: don't let the adversary pick the function.
\end{ideabox}

\section{5.8 Direct Hashing}

\emph{Direct hashing} is the degenerate case: use the key itself
as the bucket index. No modulo, no mixing, no collisions.
Whenever it's feasible, it's the fastest possible hash table.
Unfortunately, it's rarely feasible.

\begin{defnbox}[Direct access table]
A hash table in which the hash function is the identity:
$h(k) = k$. The bucket for key $k$ is always $\text{table}[k]$.
No collisions can occur because every key has its own bucket.
\end{defnbox}

\subsection*{When it works}

Direct hashing requires \emph{every} potential key to have a
reserved slot. If your keys are integers in $[0, K)$, you need a
table of size $K$.

\begin{examplebox}[Direct-access table]
\begin{lstlisting}[language=C++,basicstyle=\ttfamily\small,frame=single]
// Suitable only when keys are small non-negative integers.
constexpr size_t K = 10000;
std::optional<Item> table[K];

void insert(const Item& it)   { table[it.key] = it; }
void remove(int key)          { table[key].reset(); }
Item* find(int key) {
    return table[key] ? &*table[key] : nullptr;
}
\end{lstlisting}
\end{examplebox}

Three lines per operation, all $O(1)$ \emph{worst case} --- not
average. No hash function to compute, no collision handling, no
resizing.

\subsection*{The two killer limitations}

\begin{center}
\begin{tabular}{|l|p{9cm}|}
\hline
\textbf{Limit} & \textbf{Consequence} \\
\hline
Keys must be small non-negative integers & Can't directly handle strings, negative numbers, or large key spaces \\
Table size = max key + 1                 & A key space of $2^{32}$ needs a 4-billion-slot array (16 GB if each slot is 4 bytes, more if you store anything bigger) \\
\hline
\end{tabular}
\end{center}

If the key space is dense --- say, student IDs 0--9999 and you
have 8000 active students --- direct hashing is perfect: a
10,000-slot array, 80\% utilized, all operations $O(1)$ with zero
overhead.

If the key space is sparse --- say, US social security numbers
(9 digits, $10^9$ possible values) but your table has a few
thousand people --- direct hashing wastes 99.9999\% of the array.
You need a regular hash table.

\begin{ideabox}[Direct hashing is the ceiling]
Direct hashing represents the best a hash table can do: one
memory access per operation, zero collision work. Every
``real'' hash table is an approximation of this under memory
pressure. When you profile a hash-heavy workload and it's still
slow, the question to ask is ``how close am I to direct hashing
behavior?'' --- the answer tells you whether the slowness is
unavoidable or an artifact of your chosen scheme.
\end{ideabox}

\subsection*{Where direct hashing actually shows up}

Don't dismiss it as purely theoretical. Small-keyspace direct
access tables are genuinely common:

\begin{itemize}[leftmargin=*]
    \item \textbf{ASCII character lookup tables} (256 slots) ---
          tokenizer tables, character-class tests
          (\texttt{isalnum}, \texttt{isspace})
    \item \textbf{Byte-to-hex} conversion (256 slots holding
          two-character strings)
    \item \textbf{Color palettes} in graphics (indexed color:
          the ``hash'' is just the palette index)
    \item \textbf{File descriptor tables} in OS kernels (fd is
          already an array index)
    \item \textbf{CPU register mapping} in compilers
          (register number $\to$ allocation info)
    \item \textbf{Symbol tables in assemblers} when symbol IDs
          are sequential
\end{itemize}

Every one of these uses direct indexing by design because the
key space is small and dense.

\begin{warnbox}[Negative keys need translation]
If your keys include negative integers, you can still use direct
hashing --- just add an offset. Keys in $[-100, 100]$ map to
indices $[0, 200]$ via \texttt{index = key + 100}. This extends
direct hashing to any \emph{bounded} integer range.
\end{warnbox}

\subsection*{Comparison}

\begin{center}
\begin{tabular}{|l|c|c|c|}
\hline
\textbf{Strategy} & \textbf{Worst case} & \textbf{Collisions?} & \textbf{Space} \\
\hline
Direct hashing    & $O(1)$ & None       & $O(K)$ (key range) \\
Chaining          & $O(n)$ & Handled    & $O(n + m)$ \\
Open addressing   & $O(n)$ & Handled    & $O(m)$ \\
\hline
\end{tabular}
\end{center}

Direct hashing trades space for time: you pay for a slot for
every possible key, but you get the best possible access time.
Regular hash tables reverse the deal: save memory, accept
occasional collision handling.

\begin{notebox}[Perfect hashing: the middle ground]
For static key sets (known up front, never changes), you can
construct a \emph{perfect hash function} --- one that maps each
key to a unique bucket in a tightly-sized table. No wasted space,
no collisions. Classic use: compiler keyword tables (the set
\{\texttt{if}, \texttt{else}, \texttt{while}, \texttt{for},
\ldots\} is fixed; you can hand-tune a hash function that maps
each to a distinct small integer). Tools like \texttt{gperf}
generate perfect hash functions from a word list. It's direct
hashing for the modern world --- limited to static data but
optimal when applicable.
\end{notebox}

\begin{ideabox}[Why this section exists]
Direct hashing is deliberately a \emph{counterexample} --- it
shows what you'd do if memory were free. Every technique in
sections 5.1--5.7 exists because direct hashing's memory cost is
prohibitive. The moment you squint at a hash table problem and
realize ``oh, my keys are small integers,'' you should stop and
consider an array.
\end{ideabox}

\section{5.9 Hashing for Cryptography}

Hashing has a second life outside data structures: as a
\emph{cryptographic primitive}. Cryptographic hash functions
share the ``deterministic function from input to fixed-size
output'' shape with hash-table functions, but their design
goals --- and their cost --- are radically different.

\begin{defnbox}[Cryptographic hash function]
A hash function with three additional properties:
(1) \emph{pre-image resistance}: given a hash value $h$, it is
computationally infeasible to find any input $x$ such that
$\text{hash}(x) = h$;
(2) \emph{second pre-image resistance}: given an input $x_1$, it
is infeasible to find a different $x_2$ with the same hash;
(3) \emph{collision resistance}: it is infeasible to find any
two inputs with the same hash.
\end{defnbox}

\subsection*{What ``infeasible'' means}

A good cryptographic hash with an $n$-bit output requires
approximately $2^n$ work to find a pre-image and $2^{n/2}$ work
to find a collision (the birthday bound). For $n = 256$, that is
$\approx 10^{77}$ operations --- more than the age of the
universe on all the computers ever built.

Compare this to a hash-table hash function like djb2. djb2 is
fast (a few nanoseconds per byte) but trivially easy to reverse
or collide on demand. It's not secure; it was never designed to
be.

\begin{center}
\begin{tabular}{|l|l|l|}
\hline
\textbf{Goal} & \textbf{Hash-table hash} & \textbf{Cryptographic hash} \\
\hline
Speed          & $\leq$ 1 ns/byte    & 1--10 ns/byte (SHA-256) \\
Output size    & 32--64 bits         & 128--512 bits \\
Reversibility  & Often easy           & Infeasible \\
Collisions     & Acceptable (handled) & Infeasible to produce \\
Use cases      & Lookups              & Integrity, signatures, passwords \\
\hline
\end{tabular}
\end{center}

\subsection*{Data integrity}

Download a file plus its hash. Recompute the hash locally. If
they match, the file wasn't corrupted or tampered with.

\begin{examplebox}[Verifying a download with SHA-256]
\begin{lstlisting}[basicstyle=\ttfamily\small,frame=single]
$ sha256sum ubuntu-24.04.iso
c37f6b0e...  ubuntu-24.04.iso

# Compare against the published hash on ubuntu.com.
# Match   -> bits are correct
# Mismatch -> file is corrupt or modified
\end{lstlisting}
\end{examplebox}

Historical timeline of algorithms:
\begin{itemize}[leftmargin=*]
    \item \textbf{MD5} (1991): 128-bit output. Collision attacks
          demonstrated in 2004; \textbf{broken}, do not use for
          security.
    \item \textbf{SHA-1} (1995): 160-bit output. Collision
          demonstrated in 2017 (SHAttered); \textbf{deprecated}.
    \item \textbf{SHA-256 / SHA-3}: current standards.
          256--512 bit outputs, no known attacks faster than
          brute force.
\end{itemize}

\begin{warnbox}[Never use MD5 or SHA-1 for security]
Both have been publicly broken. They remain fine for
\emph{non-security} use (checksums, content-addressable storage
keys) where an attacker isn't trying to cause collisions. But
for anything where deception is a threat model --- signatures,
certificates, password hashing --- use SHA-256 or better.
\end{warnbox}

\subsection*{Password hashing}

When a user signs up, don't store their password. Store
\texttt{hash(password)}. At login, hash the submitted password
and compare. If the database is stolen, attackers have hashes,
not passwords --- and with a good hash, inverting them is
infeasible.

\begin{examplebox}[Naive password hashing (don't do this)]
\begin{lstlisting}[basicstyle=\ttfamily\small,frame=single]
stored = sha256(password)   // bad
\end{lstlisting}
\end{examplebox}

Two attacks on naive password hashing:

\begin{enumerate}[leftmargin=*]
    \item \textbf{Rainbow tables}: attackers precompute hashes
          for billions of common passwords. If your hash is in
          the table, password recovered instantly.
    \item \textbf{Bulk cracking}: modern GPUs compute billions of
          SHA-256 hashes per second. For short or weak passwords,
          brute force is tractable.
\end{enumerate}

The two standard defenses:

\begin{center}
\begin{tabular}{|l|p{9cm}|}
\hline
\textbf{Defense} & \textbf{What it does} \\
\hline
\textbf{Salt} & Append a random per-user value before hashing. Forces the attacker to recompute the rainbow table for every user --- no more precomputation. \\
\textbf{Slow hash} & Use a deliberately expensive function (bcrypt, scrypt, Argon2) that takes 100ms--1s per hash. Kills bulk cracking: a million-password attack now takes 12 CPU-days instead of a second. \\
\hline
\end{tabular}
\end{center}

\begin{examplebox}[Correct password hashing]
\begin{lstlisting}[basicstyle=\ttfamily\small,frame=single]
salt   = random_bytes(16)                  // unique per user
stored = bcrypt(password + salt, cost=12)  // slow
// store (salt, stored) in the database
\end{lstlisting}
\end{examplebox}

\texttt{bcrypt}, \texttt{scrypt}, and \texttt{Argon2} are the
current standards. Argon2 won the Password Hashing Competition
in 2015 and is the recommended default for new systems.

\begin{ideabox}[The design tension: fast vs. slow]
Cryptographic hashes usually want to be \emph{fast} (you hash a
lot of data for integrity). Password hashes want to be
\emph{slow} (you hash once per login, and slow is the point).
This is why ``use SHA-256 for passwords'' is wrong --- SHA-256
was engineered for speed, the opposite of what password hashing
needs. Always use a dedicated password hash for passwords.
\end{ideabox}

\subsection*{Where else does cryptographic hashing show up?}

\begin{itemize}[leftmargin=*]
    \item \textbf{Digital signatures}: sign the hash of the
          document, not the document itself. Cheaper, and works
          for any size.
    \item \textbf{Blockchain / content addressing}: each block
          references the hash of the previous block. Git does
          the same with commits.
    \item \textbf{Merkle trees}: hash-based structures that let
          you verify membership in a large set without
          downloading the whole set.
    \item \textbf{HMAC}: keyed hash functions used to authenticate
          messages.
    \item \textbf{Commitment schemes}: hash now, reveal later ---
          used in auctions, coin-flipping protocols, etc.
\end{itemize}

\begin{notebox}[Cryptographic hashes and hash tables in one program]
A modern web server might use both. SHA-256 to verify TLS
certificates and check uploaded file integrity; a fast non-
cryptographic hash (like xxHash or std::hash) to back
\texttt{unordered\_map}s in the session cache. The two systems
have different threat models and different cost profiles; it's
normal for one program to use both.
\end{notebox}

\begin{ideabox}[Chapter 5 in retrospect]
You started with a simple idea: \emph{compute the address instead
of searching for it}. Then spent eight sections on the
engineering reality: hash functions that spread keys out,
collision resolution strategies, resizing to bound load, when
direct hashing applies. This final section pivots to a completely
different use of the same mathematical tool. Cryptographic
hashing \emph{looks} like data-structure hashing but has almost
nothing else in common --- the only shared idea is ``deterministic
function, fixed-size output.'' The design space between a
4-nanosecond hash-table hash and a 100-millisecond Argon2 password
hash is enormous, and the choice depends entirely on what you're
defending against.
\end{ideabox}

\section*{What this connects to}

\begin{notebox}[Forward links to later chapters]
\begin{itemize}
    \item \textbf{Chapter 6 (BSTs).} BSTs give you a sorted alternative
    to hash tables. Choose BST (ordered set, range queries) vs.\ hash
    (fastest point lookup) -- the tradeoff is on the exam.
    \item \textbf{Graph algorithms} (opt.\ ch.~10). ``Visited'' sets in
    BFS / DFS are hash sets; Dijkstra's ``distance'' lookup is a hash
    map; hashing keeps graph traversal at expected $O(V + E)$.
    \item \textbf{Dynamic programming memoization.} Top-down DP stores
    already-computed subproblems in a hash map. The efficiency of DP
    in the average case is a hash-table claim.
    \item \textbf{C++ production idiom.}
    \texttt{std::unordered\_map<K, V>} is chaining; \texttt{std::map<K,
    V>} is a red-black tree (ch.~6 territory). Knowing which to pick
    is a real decision on every CS 300 assignment from here on.
    \item \textbf{Linked lists} from ch.~4 are what chaining buckets
    \emph{are}. The node + pointer drills from ch.~4 pay off the moment
    you trace a chaining collision.
\end{itemize}
\end{notebox}

\textbf{Companion materials.} Compact notes:
\texttt{chapters/ch\_5/notes.tex}. Practice prompts (problem $\to$
pseudocode $\to$ C++): \texttt{chapters/ch\_5/practice.md}.

\end{document}
